\chapter{The Generation Gap} \label{ch:generation-gap}
Around the middle of the twentieth century, certain linguists were developing a clear, mathematical conception of what it meant for something to be grammatical: it is generated by the grammar. And what is ``the grammar''? For these linguists, it's a set of operations together with a library that generate all and only the grammatical strings in a language. Let's set aside, for the moment, the circularity of these claims and focus on the operations and the library.

At core, there is really just a single operation, discrete combination, which turns out to be a rather unusual kind of combination in nature. Usually when we combine things, they blend or mix together in some way, like when we combine a seed, earth, water, air, and sunlight and a tree sprouts, or when we combine different metals under intense heat to create an alloy. But in this conception of grammar, discrete combination doesn't work like that. Instead, it's more like combining building blocks or Lego pieces. Each piece retains its identity and properties when it is combined with others, but the resulting structure has its own properties that are more than the sum of its parts.

To give an example, consider the words \data{a}, \data{dust}, and \data{storm}. Independently, these words have unique meanings and properties. \data{A} is an article, and it delimits one item or individual. \data{Dust} is a noun, referring to fine, dry particles of matter. And \data{storm} is another noun, which denotes a violent disturbance of the atmosphere with strong winds and usually rain, thunder, lightning, or snow.

Because we're using discrete combination, we couldn't, for instance, first take these words apart letter by letter and then re-combine them into \data{darts at sumo} or \data{dams out rats}. Each word is discrete. But when we apply the operation of discrete combination, we generate the phrase \data{A dust storm}, which preserves the individual words. This new structure has a distinct meaning and properties that result from the specific combination of its components. In this case, it's a single event involving high winds and dust. Each individual word retains its own meaning, but the combination conveys something different.

It also has different properties from its constituents. For example, \data{dust} and \data{storm} each have the ability to be combined with \data{cosmic}, yielding \data{cosmic dust} or \data{cosmic storm}, but if we try to combine the discrete elements \data{cosmic} and \data{a dust storm}, the resulting *\data{cosmic a dust storm} is ungrammatical. And even though \data{a cosmic dust storm} is grammatical, we are treating \data{a dust storm} as a discrete unit, so \data{cosmic} can't intrude. Instead, to get that, we'd need to first create \data{dust storm}, then combine with \data{cosmic}, and only then combine with \data{a}. 

This shows us that, in this view of things, we can't disregard the order in which combination operations are performed. Nor can we just combine things going from left to right. If we tried that, our first constituent would be *\data{a cosmic}, and that's ungrammatical.

Discrete combination is a very powerful and interesting idea. But it depends on the second element, the ``library''---the set of all words and morphemes that the language's speakers know and use. This library includes not just the words' meanings, but also information about how they can be combined with other words. For example, \data{cosmic} can be combined with \data{dust} to form \data{cosmic dust}, but not with \data{a dust storm} to form *\data{cosmic a dust storm}.

But if we are working at this level of detail, then our grammar simply becomes a list of all the possible combinations of words and phrases. And to a certain extent, this may be the case. \citet[47]{Culicover1999}

\ea
    \ea[]{\textit{It is likely that Robin will be elected president next year.}}
    \ex[]{\textit{It is probable that Robin will be elected president next year.}}
    \ex[]{\textit{Robin is likely to be elected president next year.}}
    \ex[*]{\textit{Robin is probable to be elected president next year.}}
    \z
\z


But this seems highly inefficient. Wouldn't such a grammar present significant challenges in storage and use? One problem would arise when we hear something new. Every unfamiliar combination would either have to be rejected, meaning we couldn't learn, or it would have to be incorporated it into our grammar. In such a world, we would have no concept at all of which combinations are grammatical and which are not. This suggests that there's more to grammar than just a list of possible word combinations.

It seems much better to employ some kind of abstraction that would allow more generalization. For example, since \data{dust} and \data{storm} are both nouns, and \data{cosmic} can combine with each of them, perhaps we could capture this simply with a rule that says \data{cosmic} may be combined with a noun. We could then remove specific reference to \data{cosmic} and say simply that an adjective can combine with a noun. Add a rule that an article (like \data{a}) can combine with this, and away we go. This combination could combine with a verb like \data{shimmered} to produce the full sentence: \data{A cosmic storm shimmered}.

What made these systems really clever was realizations such as that we could combine multiple Adjectives but not, for instance, multiple articles. So whatever kind of operation combined an adjective and a noun (\data{cosmic storm}) could be employed with the resulting combination (\data{giant cosmic storm}), and so on (\data{new giant cosmic storm}), and so on (\data{beautiful new giant cosmic storm}). Let's say we call an adjective + noun combination a ``nominal''. Then we can capture this recursive procedure with a rule that says an adjective plus a nominal can combine to form a nominal. And then we can combine an article with a nominal to form a different kind of structure. We'll call it a ``noun phrase''. In this way, our grammar can generate a huge number of different noun phrases with just a few simple rules.

This means that the library includes abstractions: noun, nominal, and noun phrase, along with adjective, article, and sentence. If the linguists are able to articulate all the rules about which operations are allowed, then those, together with the library should produce all and only the grammatical constituents in a language. At least, this was the way most linguists were thinking about it in the late sixties.

By the early seventies, this broad consensus had crumbled, with a number of Chomsky's former students declaring that the whole idea of grammaticality was incoherent. \citep{Harris2021}

Jim McCawley told the \textit{New Yorker}, ``Chomsky assumes that there are sentences which belong to the language and other sequences of words which don't—and the grammarian's task is to write rules [that] determine which belong and which don't. Postal and Lakoff and I say this isn't a coherent notion'' \citep{Shenker1972}.

So how does the system overcome the \textit{likely}/\textit{probable} problem that Culicover pointed out? By not being static. The way the language evolves is towards systematicity, while at the same time allowing new irregularities into the system in a constant yin and yang of order and disorder.

\section{Generational failures}
Boeckx \citep[79]{Boeckx2012} pointed out that island repair shouldn't happen in
derivational approaches that claim the grammar generates only grammatical language. It is ``a very serious puzzle for such approaches if it indeed turns out that the relevant extraction can be shown to have taken place, only to be legitimized by a subsequent, post-syntactic operation. How can something be repaired if it was never constructed in the first place? How can something be legitimized in certain contexts if it is predicted never to happen in any context?''